\section*{अध्याय \ref{ch:b2:ligand-field} --- लिगैंड क्षेत्र सिद्धांत और रंग}

\begin{solution}{exo:b2:ligand-field:1}
$d^3$: $t_{2g}^3$, CFSE $-1.2\Delta_o$। $d^8$: $t_{2g}^6e_g^2$, CFSE $-2.4 + 1.2 = -1.2\Delta_o$।
\end{solution}

\begin{solution}{exo:b2:ligand-field:2}
\ce{[Fe(H2O)6]^2+}, $d^6$ \omterm{def:b2:ligand-field:spin}{उच्च चक्रण} $t_{2g}^4e_g^2$: चार। \ce{[Fe(CN)6]^4-}, \omterm{def:b2:ligand-field:spin}{निम्न चक्रण}
$t_{2g}^6$: एक भी नहीं (\omterm{def:b2:diatomic-mos:magnetism}{प्रतिचुंबकीय})।
\end{solution}

\begin{solution}{exo:b2:ligand-field:3}
600 nm नारंगी प्रकाश है; संकुल नीला (हरिताभ नीला) दिखता है।
\end{solution}

\begin{solution}{exo:b2:ligand-field:4}
$\ce{Cl-} < \ce{H2O} < \ce{NH3} < \ce{CN-}$.
\end{solution}

\begin{solution}{exo:b2:ligand-field:5}
जल: $\Delta_o < P$, \omterm{def:b2:ligand-field:spin}{उच्च चक्रण}, $t_{2g}^3e_g^2$, पाँच अयुग्मित इलेक्ट्रॉन। सायनाइड: $\Delta_o > P$,
\omterm{def:b2:ligand-field:spin}{निम्न चक्रण}, $t_{2g}^5$, एक अयुग्मित इलेक्ट्रॉन।
\end{solution}

\begin{solution}{exo:b2:ligand-field:6}
चतुष्फलकीय \ce{[CoCl4]^2-} में विपाटन किसी अष्टफलकीय विपाटन का लगभग $4/9$ है, और
क्लोराइड जल से दुर्बल लिगैंड है: d–d पट्टी लंबी तरंगदैर्ध्यों की ओर खिसकती है (नारंगी
से लाल), और संकुल नीला दिखता है। चतुष्फलकीय संकुल अधिक प्रबल अवशोषण भी करते हैं,
इसलिए रंग गहरा है।
\end{solution}

\begin{solution}{exo:b2:ligand-field:7}
\ce{Ni^2+} $d^8$ है। वर्ग समतलीय: चार निचले कक्षक आठों इलेक्ट्रॉनों को युग्मों में रखते हैं,
$\mathrm d_{x^2-y^2}$ खाली रहता है: \omterm{def:b2:diatomic-mos:magnetism}{प्रतिचुंबकीय}। चतुष्फलकीय: $e^4t_2^4$, दो अयुग्मित इलेक्ट्रॉन
$t_2$ समुच्चय में।
\end{solution}

\begin{solution}{exo:b2:ligand-field:8}
$\tilde\nu = 1/(500 \times 10^{-7}\,\text{cm}) = \qty{20000}{cm^{-1}}$; $\Delta_o =
0.1196/(500 \times 10^{-9}) = \qty{2.39e5}{J/mol}$, अर्थात् \qty{239}{kJ/mol}।
\end{solution}

\begin{solution}{exo:b2:ligand-field:9}
\ce{Zn^2+} $d^{10}$ है और \ce{Sc^3+} $d^0$: कोई \omterm{def:b2:ligand-field:d-d}{d–d संक्रमण} संभव नहीं (\ce{Zn^2+} के लिए
कोई खाली d स्तर नहीं, \ce{Sc^3+} के लिए कोई d इलेक्ट्रॉन नहीं), इसलिए दृश्य क्षेत्र में कोई अवशोषण नहीं।
\end{solution}

\begin{solution}{exo:b2:ligand-field:10}
Co(III) $d^6$ + अमोनिया के छह एकाकी युग्म: 18 इलेक्ट्रॉन। बारह छहों \omterm{def:b2:diatomic-mos:bonding}{आबंधी कक्षकों} को भरते हैं,
छह अनाबंधी $t_{2g}$ को (अमोनिया इतना प्रबल क्षेत्र देती है कि कोबाल्ट(III) के साथ
\omterm{def:b2:ligand-field:spin}{निम्न चक्रण} हो); $e_g^*$ खाली है। सभी इलेक्ट्रॉन युग्मित हैं: \omterm{def:b2:diatomic-mos:magnetism}{प्रतिचुंबकीय}।
\end{solution}

\begin{solution}{exo:b2:ligand-field:11}
\ce{CO} एक अच्छा $\sigma$ दाता है और, सबसे बढ़कर, एक $\pi$ ग्राही: उसके खाली $\pi^*$ कक्षक
\omterm{def:b2:ligand-field:pi-ligands}{पश्च-दान} द्वारा $t_{2g}$ स्तर को नीचे लाते हैं, जिससे $\Delta_o$ चौड़ा होता है। \ce{F-} एक $\pi$ दाता है, जो
$t_{2g}$ को ऊपर उठाता है और $\Delta_o$ को संकरा करता है। बिंदु-आवेश मॉडल केवल आवेश देखता है और
क्रम ग़लत पाता है।
\end{solution}

\begin{solution}{exo:b2:ligand-field:12}
क्षेत्र के बिना जलयोजन एन्थैल्पियाँ एक सहज रेखा का अनुसरण करतीं (श्रेणी के अनुदिश आयन
सिकुड़ते हैं)। जल उच्च-चक्रण ऐक्वा संकुल देता है जिनकी CFSE $d^0$, $d^5$, $d^{10}$ के लिए शून्य
और $d^3$ तथा $d^8$ के लिए सबसे बड़ी है: अतिरिक्त स्थायीकरण दो कूबड़ जोड़ता है, जिनके शिखर
\ce{V^2+} और \ce{Ni^2+} के निकट हैं, उस रेखा के ऊपर जो \ce{Ca^2+}, \ce{Mn^2+}, \ce{Zn^2+} से होकर जाती है।
\end{solution}

\begin{solution}{pb:b2:ligand-field:1}
\textbf{1.} $d^3$।
\textbf{2.} $t_{2g}^3$, तीनों कक्षकों में से हर एक में एक इलेक्ट्रॉन, समांतर चक्रण।
\textbf{3.} नहीं: तीन इलेक्ट्रॉन बिना युग्मन के $t_{2g}$ में समा जाते हैं।
\textbf{4.} $-1.2\Delta_o$।
\textbf{5.} तीन।
\textbf{6.} CFSE बड़ी है और $e_g^*$ कक्षक, जो लिगैंडों की ओर इंगित करते हैं,
खाली हैं: कोई लिगैंड हटाने या जोड़ने में उस स्थायीकरण का बड़ा भाग खर्च होता है।
\textbf{7.} हल्के विकृत अष्टफलक के कोनों पर छह ऑक्साइड आयन।
\textbf{8.} $1/(556 \times 10^{-7}) = \qty{17990}{cm^{-1}}$, लगभग \qty{1.80e4}{cm^{-1}}।
\textbf{9.} $0.1196/(556 \times 10^{-9}) = \qty{2.15e5}{J/mol}$: \qty{215}{kJ/mol}।
\textbf{10.} पीला-हरा (556 nm) और बैंगनी (405 nm)।
\textbf{11.} लाल, थोड़े नीले के साथ: माणिक का गहरा लाल।
\textbf{12.} किसी $d^3$ आयन के एक से अधिक उत्तेजित विन्यास होते हैं जिनमें एक इलेक्ट्रॉन $e_g$ में हो;
इलेक्ट्रॉन प्रतिकर्षण उन्हें भिन्न ऊर्जाएँ देता है, इसलिए कई पट्टियाँ (जिनका विवेचन
तृतीय वर्ष के खंड में है)।
\textbf{13.} $0.1196/(620 \times 10^{-9}) = \qty{1.93e5}{J/mol}$, \qty{193}{kJ/mol}
(\qty{16130}{cm^{-1}})।
\textbf{14.} माणिक: उसका विपाटन बड़ा है।
\textbf{15.} नारंगी-लाल की ओर: अब लाल अवशोषित होता है, और हरा (कुछ नीले के साथ) पारगमित होता है।
\textbf{16.} बेरिल में क्रोमियम स्थल थोड़े बड़े हैं और ऑक्साइड आयन अधिक दूर,
और परिवेश भिन्न है: दुर्बल क्षेत्र।
\textbf{17.} नहीं: $d^3$ आयन किसी भी अष्टफलकीय क्षेत्र में तीन अयुग्मित इलेक्ट्रॉन रखता है।
\textbf{18.} $d^9$।
\textbf{19.} ऐक्वा संकुल \omterm{def:b2:ligand-field:d-d}{d–d संक्रमण} द्वारा लाल और नारंगी में अवशोषण करता है:
पारगमित प्रकाश नीला है।
\textbf{20.} यह सफ़ेद हो जाता है: कॉपर के चारों ओर जल लिगैंडों के बिना (सल्फ़ेट उनका
स्थान लेता है) क्षेत्र दुर्बल होता है और पट्टी अवरक्त में चली जाती है।
\textbf{21.} जल फिर कॉपर से बँधता है और नीला ऐक्वा संकुल बनता है।
\textbf{22.} छोटी तरंगदैर्ध्य पर: श्रेणी में जल से ऊपर स्थित अमोनिया बड़ा
विपाटन देती है (ऐम्मीन संकुल का गहरा नीला-बैंगनी)।
\textbf{23.} नहीं: $t_{2g}^6e_g^3$ एकमात्र विन्यास है।
\textbf{24.} हाँ: एक अयुग्मित इलेक्ट्रॉन।
\textbf{25.} $\boldsymbol{\Delta_o \approx \qty{1.8e4}{cm^{-1}}}$, लगभग
$\boldsymbol{\qty{215}{kJ/mol}}$।
\end{solution}
